\section*{Chapitre \ref{ch:b1:potential-capacitors} --- Potentiel électrique~; conducteurs et condensateurs}

\begin{solution}{exo:b1:potential-capacitors:1}
$V = 8.99 \times 10^9 \times 10^{-6}/1 = \qty{9.0}{kV}$~; $W = qV = \qty{9.0}{mJ}$~;
libérée, $\tfrac12mv^2 = \qty{9}{mJ}$~: $v = \sqrt{18} = \qty{4.2}{m/s}$.
\end{solution}

\begin{solution}{exo:b1:potential-capacitors:2}
$E = 100/10^{-3} = \qty{1e5}{V/m}$~; $E_k = eU = \qty{100}{eV} = \qty{1.6e-17}{J}$,
donc $v = \sqrt{2E_k/m_e} = \qty{5.9e6}{m/s}$~; le mouvement étant uniformément
accéléré, $t = 2d/v = \qty{3.4e-10}{s}$.
\end{solution}

\begin{solution}{exo:b1:potential-capacitors:3}
$E_A = 10/0.002 = \qty{5}{kV/m}$, $E_B = 10/0.008 = \qty{1.25}{kV/m}$. Le champ
est dirigé vers les potentiels décroissants~; un électron ($q < 0$) accélère en
sens inverse, vers les $V$ croissants.
\end{solution}

\begin{solution}{exo:b1:potential-capacitors:4}
$C = 4\pi\varepsilon_0R = 1.11 \times 10^{-10} \times 0.15 = \qty{17}{pF}$~; $Q = CV =
\qty{3.3}{\micro C}$~; $E = V/R = 2 \times 10^5/0.15 = \qty{1.3}{MV/m}$. Claquage
à $V = E_{\max}R = 3 \times 10^6 \times 0.15 = \qty{450}{kV}$ --- les grosses
machines emploient de grosses sphères.
\end{solution}

\begin{solution}{exo:b1:potential-capacitors:5}
Tout élément de l'anneau est à la distance $\sqrt{R^2 + z^2}$~: $V =
Q/4\pi\varepsilon_0\sqrt{R^2 + z^2}$~; $E_z = -\dd V/\dd z = Qz/4\pi\varepsilon_0(R^2 +
z^2)^{3/2}$. Disque~: $\dd q = 2\pi\sigma a\,\dd a$ à la distance $\sqrt{a^2 + z^2}$~: $V =
\dfrac{\sigma}{2\varepsilon_0}\displaystyle\int_0^R\dfrac{a\,\dd a}{\sqrt{a^2 + z^2}} =
\dfrac{\sigma}{2\varepsilon_0}\bigl(\sqrt{z^2 + R^2} - z\bigr)$, et $-\dd V/\dd z =
\dfrac{\sigma}{2\varepsilon_0}\Bigl(1 - \dfrac{z}{\sqrt{z^2 + R^2}}\Bigr)$.
\end{solution}

\begin{solution}{exo:b1:potential-capacitors:6}
$V(r) = V(R) + \displaystyle\int_r^R\dfrac{Qr'\,\dd r'}{4\pi\varepsilon_0R^3} =
\dfrac{Q}{4\pi\varepsilon_0R} + \dfrac{Q(R^2 - r^2)}{8\pi\varepsilon_0R^3} = \dfrac{Q}{4\pi\varepsilon_0}
\dfrac{3R^2 - r^2}{2R^3}$. Pour l'or~: $V(0) = 1.5 \times 8.99 \times 10^9 \times 1.26 \times
10^{-17}/7 \times 10^{-15} = \qty{24}{MV}$~: il faut \qty{24}{MeV} à un proton et
\qty{48}{MeV} à une particule alpha ($2e$) pour atteindre le centre, \qty{32}{MeV}
pour toucher la surface. Les alphas de \qty{5}{MeV} de Rutherford rebroussent chemin en $r$ tel que $2eV =
\qty{5}{MeV}$~: $V = \qty{2.5}{MV}$, $r = \qty{45}{fm}$ --- bien à l'extérieur
du noyau, ce qui explique que son analyse en charge ponctuelle ait marché.
\end{solution}

\begin{solution}{exo:b1:potential-capacitors:7}
Sur l'axe~: $2p/4\pi\varepsilon_0r^3 = 2 \times 8.99 \times 10^9 \times 6.2 \times 10^{-30}/10^{-27} =
\qty{1.1e8}{V/m}$~; perpendiculairement, la moitié~: \qty{5.6e7}{V/m}. Moment $pE =
\qty{6.2e-24}{N.m}$~; écart d'énergie $2pE = \qty{1.2e-23}{J}$ contre $k_BT =
\qty{4.1e-21}{J}$~: $330$ fois moins --- seulement $0.1\%$ d'alignement moyen.
Ion sodium à \qty{0.3}{nm}~: $E = e/4\pi\varepsilon_0r^2 = \qty{1.6e10}{V/m}$, $-pE
= \qty{-1.0e-19}{J} = \qty{-0.6}{eV}$, soit $24\,k_BT$~: les molécules d'eau se
verrouillent sur l'ion --- c'est l'hydratation.
\end{solution}

\begin{solution}{exo:b1:potential-capacitors:8}
$Q = 4\pi\varepsilon_0R^2E = 10^{-4} \times 3 \times 10^6/8.99 \times 10^9 = \qty{33}{nC}$,
$V = ER = \qty{30}{kV}$. Sphères reliées~: $E_1 = V/R_1 = \qty{0.2}{MV/m}$,
$E_2 = V/R_2 = \qty{20}{MV/m} > \qty{3}{MV/m}$~: la petite sphère ionise l'air
et laisse fuir la charge --- c'est l'\omterm{ex:b1:potential-capacitors:point}{effet de pointe}.
\end{solution}

\begin{solution}{exo:b1:potential-capacitors:9}
$C = \varepsilon_r\varepsilon_0S/d = 2.2 \times 8.85 \times 10^{-12}/10^{-5} = \qty{2.0}{\micro F}$~;
l'enroulement change la forme, non l'aire en regard~: même $C$ (avec un second
film, de sorte que chaque feuille fasse face à l'autre des deux côtés, il double).
$W = \tfrac12 \times 2 \times 10^{-6} \times 400^2 = \qty{0.16}{J}$. Défibrillateur~:
$Q = 150 \times 10^{-6} \times 2000 = \qty{0.30}{C}$, $W = \tfrac12CU^2 = \qty{300}{J}$,
$P = 300/0.005 = \qty{60}{kW}$.
\end{solution}

\begin{solution}{exo:b1:potential-capacitors:10}
(a) $E = Q/4\pi\varepsilon_0r^2$, $U = \dfrac{Q}{4\pi\varepsilon_0}\Bigl(\dfrac1{R_1} -
\dfrac1{R_2}\Bigr)$, $C = 4\pi\varepsilon_0R_1R_2/(R_2 - R_1)$~; avec $R_2 - R_1 = d \ll
R_1$, $R_1R_2 \approx R^2$ et $C \approx 4\pi R^2\varepsilon_0/d = \varepsilon_0S/d$~; avec $R_2 \to
\infty$, $C \to 4\pi\varepsilon_0R_1$. (b) $C/L = 2\pi\varepsilon_r\varepsilon_0/\ln(b/a) = 2\pi \times 2.3
\times 8.85 \times 10^{-12}/\ln5 = \qty{79}{pF/m}$, le bon ordre. (c) $E(a) =
U/[a\ln(b/a)] = 1000/(5 \times 10^{-4} \times 1.61) = \qty{1.2}{MV/m}$~: au-dessous
des \qty{20}{MV/m} du polyéthylène, sans danger~; dans l'air ce serait limite.
\end{solution}

\begin{solution}{exo:b1:potential-capacitors:11}
(a) $W = Q^2/2C = Q^2/8\pi\varepsilon_0R$. (b) $R = e^2/8\pi\varepsilon_0m_ec^2 = 2.31 \times
10^{-28}/(2 \times 8.19 \times 10^{-14}) = \qty{1.4e-15}{m}$ (le «~rayon classique
de l'électron~» usuel, $e^2/4\pi\varepsilon_0m_ec^2 = \qty{2.8}{fm}$, omet le $\tfrac12$). (c) $W =
Q^2/2C = Q^2d/2\varepsilon_0S$ double~: $\Delta W = Q^2d/2\varepsilon_0S = F\,d$ avec $F =
Q^2/2\varepsilon_0S = \sigma^2S/2\varepsilon_0$, comme trouvé directement. (d) À $U$ fixée, $W =
\varepsilon_0SU^2/2d$ est divisée par deux~: $\Delta W = -W/2$~; celui qui tire fournit
encore $W/2$ (en intégrant $F = \varepsilon_0SU^2/2x^2$ de $d$ à $2d$)~; les deux
moitiés, soit $W$ en tout, repartent dans la pile lorsque la charge $CU/2$ lui
revient sous la tension $U$.
\end{solution}

\begin{solution}{exo:b1:potential-capacitors:12}
(a) La sphère de rayon $r$ porte $q = Qr^3/R^3$ et son potentiel de surface vaut
$q/4\pi\varepsilon_0r$~; ajouter la coquille $\dd q = 3Qr^2\,\dd r/R^3$ coûte $\dd W =
q\,\dd q/4\pi\varepsilon_0r = 3Q^2r^4\,\dd r/4\pi\varepsilon_0R^6$~; on intègre~: $W =
3Q^2/20\pi\varepsilon_0R$. (b) $Q = 92e = \qty{1.47e-17}{C}$~: $W = 3 \times 2.17 \times
10^{-34}/(20\pi \times 8.85 \times 10^{-12} \times 7.4 \times 10^{-15}) = \qty{1.6e-10}{J}
= \qty{990}{MeV}$. (c) Chaque fragment~: charge $Q/2$, rayon $R/2^{1/3}$,
énergie $\tfrac14 2^{1/3}W$~; les deux ensemble~: $2^{1/3}W/2 = 0.63W$~: on libère
$0.37 \times 990 = \qty{370}{MeV}$. (d) La «~tension superficielle~» nucléaire~:
deux fragments ont plus de surface qu'un noyau, ce qui coûte environ
\qty{170}{MeV}~; le solde, voisin de \qty{200}{MeV}, apparaît surtout en \omterm{thm:b1:work-and-energy:ke}{énergie
cinétique} des deux fragments qui s'éloignent sous leur répulsion
mutuelle.
\end{solution}

\begin{solution}{pb:b1:potential-capacitors:1}
\textbf{1.} Le poids $\tfrac43\pi r^3\rho g$ vers le bas, la traînée $6\pi\eta rv$
vers le haut, la \omterm{thm:b1:fluid-statics:archimedes}{poussée d'Archimède} $\tfrac43\pi r^3\rho_{\text{air}}g$ vers le
haut --- soit $1.2/900 = 0.13\%$ du poids, négligeable. Régime limite~: $\tfrac43\pi r^3\rho g = 6\pi\eta rv_1$, $v_1 =
2\rho gr^2/9\eta$.

\textbf{2.} $v_1 = \qty{1.0e-4}{m/s}$~: $r = \sqrt{9\eta v_1/2\rho g} = \sqrt{9 \times
1.8 \times 10^{-5} \times 10^{-4}/(2 \times 900 \times 9.8)} = \qty{9.6e-7}{m}$, un
micromètre~; $m = \tfrac43\pi r^3\rho = \qty{3.3e-15}{kg}$.

\textbf{3.} $\tau = m/6\pi\eta r = 3.3 \times 10^{-15}/3.3 \times 10^{-10} = \qty{1e-5}{s}$~;
distance $\sim v_1\tau = \qty{1}{nm}$. Jamais.

\textbf{4.} $\mathrm{Re} = 1.2 \times 10^{-4} \times 9.6 \times 10^{-7}/1.8 \times 10^{-5} =
6 \times 10^{-6} \ll 1$~: la \omterm{prop:b1:newton-dynamics:drag}{loi de Stokes} est valable.

\textbf{5.} $eE = 1.6 \times 10^{-19} \times 3 \times 10^5 = \qty{4.8e-14}{N}$ contre
$mg = \qty{3.2e-14}{N}$~: comparables seulement pour des gouttes microniques. Une
goutte de \qty{10}{\micro m} pèse mille fois plus~; une \omterm{def:b1:electrostatics-gauss:charge}{charge élémentaire} ne
changerait sa vitesse que de $0.1\%$, invisible.

\textbf{6.} Le mouvement brownien~: le bombardement par les molécules d'air
secoue la goutte au hasard~; les chronométrages se dispersent, ce qui limite la
précision et impose une taille minimale de goutte utile.

\textbf{7.} $E = U/d = 5000/0.016 = \qty{3.1e5}{V/m}$~; les armatures sont
beaucoup plus larges que $d$, de sorte que la superposition des deux plans donne
un champ uniforme~; les équipotentielles sont des plans horizontaux, $V = Uz/d$.

\textbf{8.} $C = \varepsilon_0\pi(0.1)^2/0.016 = \qty{17}{pF}$~; $Q = CU = \qty{87}{nC}$~;
$W = \tfrac12CU^2 = \qty{0.22}{mJ}$.

\textbf{9.} L'armature supérieure étant positive, $\vect E$ est dirigé vers le
bas~; la force devant être dirigée vers le haut, $q < 0$. $|q| = mgd/U_0 = 3.3 \times 10^{-15} \times 9.8 \times 0.016/
1080 = \qty{4.8e-19}{C} = 3e$~: $q = -3e$.

\textbf{10.} En montée à vitesse constante~: $|q|E = mg + 6\pi\eta rv_2$ et $mg =
6\pi\eta rv_1$, donc $|q|E = 6\pi\eta r(v_1 + v_2)$ et $|q| = 6\pi\eta r(v_1 +
v_2)d/U$. $v_2 = (|q|E - mg)/6\pi\eta r = (1.5 \times 10^{-13} - 3.2 \times 10^{-14})/
3.3 \times 10^{-10} = \qty{3.6e-4}{m/s}$~: \qty{1}{mm} en \qty{2.8}{s}.

\textbf{11.} L'équilibre est un jugement de zéro porté sur un grain qui dérive
et tremble, et dont la charge peut changer entre-temps~; les chronométrages sur
une distance fixée se répètent à volonté, $r$ est fixé une fois pour toutes par
$v_1$, et une seule goutte fournit alors toute une série de charges.

\textbf{12.} $|q|U = 3 \times 1.6 \times 10^{-19} \times 5000 = \qty{2.4e-15}{J}$
(\qty{15}{keV}) sur tout l'intervalle~; traînée sur \qty{1}{cm}~: $(|q|E - mg) \times 0.01 =
\qty{1.2e-15}{J}$, convertis en chaleur dans l'air.

\textbf{13.} En divisant par $1.60$~: $3.00$, $5.01$, $4.01$, $1.99$, $6.02$ ---
soit les entiers $3, 5, 4, 2, 6$.

\textbf{14.} $q/n$~: $1.600, 1.602, 1.605, 1.595, 1.605$~; moyenne $1.601$,
écart-type $0.004$, incertitude sur la moyenne $0.004/\sqrt5 =
0.002$~: $e = (1.601 \pm 0.002) \times 10^{-19}$~C.

\textbf{15.} Les rayons X ionisent l'air~; la goutte capte les ions un à un,
chacun portant un nombre entier de $e$.

\textbf{16.} $r \propto \eta^{1/2}$ d'après $v_1$, $m \propto r^3 \propto \eta^{3/2}$, et
$|q| = mgd/U_0 \propto \eta^{3/2}$~: $1\%$ sur $\eta$ donne $1.5\%$ sur $e$. Le
$\eta$ de Millikan était $0.4\%$ trop bas, d'où son $e$ $0.6\%$ trop bas.

\textbf{17.} Les quarks ne viennent jamais seuls~: ils sont confinés dans les
hadrons, et tout objet libre porte une charge entière.

\textbf{18.} $N_A = F/e = 96485/1.602 \times 10^{-19} = \qty{6.02e23}{mol^{-1}}$.

\textbf{19.} $m_e = 1.602 \times 10^{-19}/1.76 \times 10^{11} = \qty{9.1e-31}{kg}$~: $1840$
fois plus léger que l'atome d'hydrogène.

\textbf{20.} Des plans horizontaux entre les armatures, qui se bombent près des
bords~; des \omterm{def:b1:electrostatics-gauss:field}{lignes de champ} verticales, s'évasant sur les bords sur une largeur
de l'ordre de $d = \qty{1.6}{cm}$ --- peu de chose devant les \qty{20}{cm} de
diamètre, et l'on observe la goutte au centre.

\textbf{21.} $E_p = qUz/d + mgz$, linéaire en $z$~: pas de minimum~; une goutte
immobile exige que la pente s'annule, $q = -mgd/U$, et l'équilibre est alors
indifférent --- toute dérive résiduelle persiste, ce qui rend l'équilibre
difficile à juger.

\textbf{22.} $2.4 \times 10^{-15}$ contre $2.2 \times 10^{-4}$~J~: $10^{-11}$. La goutte
ne perturbe pas le champ.

\textbf{23.} Champ vers le haut, force vers le bas sur la goutte négative~: elle tombe à $(mg
+ |q|E)/6\pi\eta r = (3.2 \times 10^{-14} + 1.5 \times 10^{-13})/3.3 \times 10^{-10} =
\qty{5.6e-4}{m/s}$.

\textbf{24.} $F = 10^4 \times 1.6 \times 10^{-19} \times 3.1 \times 10^5 = \qty{5e-10}{N}$
contre $\qty{9.8e-9}{N}$~: $5\%$ du poids --- pas de lévitation possible, et avec
$10^4$ charges un changement d'un $e$ serait invisible. Un brouillard d'huile
donne des gouttes petites, sphériques et qui ne s'évaporent pas.

\textbf{25.} Mesurés~: $v_1$, $v_2$ et $U$ (avec $d$, $\eta$, $\rho$ connus)~;
déduits~: $r$, $m$, puis $q$, toujours multiple entier de
$e = (1.601 \pm 0.002) \times 10^{-19}$~C ($0.1\%$)~; il en découle $N_A = F/e$ et
$m_e = e/(e/m_e)$.
\end{solution}
